\section{Segment Tree Animations}

\subsection{Standard Segment Tree -- Query and Update}

Build a segment tree from $a = [2, 5, 1, 3, 7, 4]$ with sum aggregation.
First we query the range $[1, 4]$, then perform a point update at index 2.

\begin{animation}[id="segtree-query-update", label="Segment tree query and point update"]
\shape{st}{Tree}{data=[2,5,1,3,7,4], kind="segtree", show_sum=true}
\shape{arr}{Array}{size=6, data=[2,5,1,3,7,4], labels="0..5", label="$a$"}

\step
\narrate{Segment tree built from $a = [2, 5, 1, 3, 7, 4]$. Each internal node stores the sum of its range. Tree nodes: $[0,5] \to [0,2],[3,5] \to [0,1],[2,2],[3,4],[5,5] \to [0,0],[1,1],[3,3],[4,4]$.}

\step
\highlight{arr.range[1:4]}
\narrate{Query: compute $\text{sum}(1, 4)$. The query range $[1, 4]$ is highlighted in the source array.}

\step
\recolor{st.node["[0,5]"]}{state=current}
\narrate{Start at root $[0,5]$. The query $[1,4]$ partially overlaps, so we recurse into both children.}

\step
\recolor{st.node["[0,2]"]}{state=current}
\recolor{st.node["[0,0]"]}{state=dim}
\narrate{Left subtree $[0,2]$: recurse down. Node $[0,0]$ is outside query range -- skip it.}

\step
\recolor{st.node["[0,1]"]}{state=current}
\recolor{st.node["[1,1]"]}{state=good}
\recolor{st.node["[2,2]"]}{state=good}
\narrate{Node $[1,1]$ fully inside $[1,4]$ -- take value $5$. Node $[2,2]$ fully inside -- take value $1$.}

\step
\recolor{st.node["[3,5]"]}{state=current}
\recolor{st.node["[3,4]"]}{state=good}
\recolor{st.node["[5,5]"]}{state=dim}
\narrate{Right subtree $[3,5]$: node $[3,4]$ fully covered -- take sum $= 10$. Node $[5,5]$ outside range -- skip.}

\step
\highlight{st.node["[1,1]"]}
\highlight{st.node["[2,2]"]}
\highlight{st.node["[3,4]"]}
\annotate{st.node["[0,5]"]}{label="ans = 16", color=good}
\narrate{Result: $5 + 1 + 10 = 16$. Three nodes contributed -- this is why queries are $O(\log n)$.}

\step
\recolor{st.node["[0,5]"]}{state=idle}
\recolor{st.node["[0,2]"]}{state=idle}
\recolor{st.node["[3,5]"]}{state=idle}
\recolor{st.node["[0,0]"]}{state=idle}
\recolor{st.node["[0,1]"]}{state=idle}
\recolor{st.node["[1,1]"]}{state=idle}
\recolor{st.node["[2,2]"]}{state=idle}
\recolor{st.node["[3,4]"]}{state=idle}
\recolor{st.node["[3,3]"]}{state=idle}
\recolor{st.node["[4,4]"]}{state=idle}
\recolor{st.node["[5,5]"]}{state=idle}
\narrate{Now perform a point update: set $a[2] = 10$ (was 1). We walk from the leaf up to root.}

\step
\recolor{arr.cell[2]}{state=current}
\apply{arr.cell[2]}{value=10}
\recolor{st.node["[2,2]"]}{state=current}
\apply{st.node["[2,2]"]}{value=10}
\narrate{Update leaf $[2,2]$: value changes from 1 to 10.}

\step
\recolor{st.node["[2,2]"]}{state=done}
\recolor{st.node["[0,2]"]}{state=current}
\apply{st.node["[0,2]"]}{value=17}
\narrate{Update $[0,2]$: new sum $= [0,1]\text{.sum} + 10 = 7 + 10 = 17$.}

\step
\recolor{st.node["[0,2]"]}{state=done}
\recolor{st.node["[0,5]"]}{state=current}
\apply{st.node["[0,5]"]}{value=31}
\narrate{Update root $[0,5]$: new sum $= 17 + 14 = 31$. Point update complete in $O(\log n)$.}

\step
\recolor{st.node["[0,5]"]}{state=done}
\recolor{arr.cell[2]}{state=done}
\narrate{Final tree reflects the updated array $a = [2, 5, 10, 3, 7, 4]$ with root sum $= 31$.}
\end{animation}


\subsection{Sparse Segment Tree -- Dynamic Node Allocation}

A sparse segment tree over range $[0, 7]$. Nodes are allocated on-demand as values are inserted.

\begin{animation}[id="sparse-segtree-demo", label="Sparse segment tree dynamic allocation"]
\shape{sst}{Tree}{kind="sparse_segtree", range_lo=0, range_hi=7}

\step
\narrate{Sparse segment tree on range $[0, 7]$. Initially only the root $[0,7]$ exists. All other nodes are virtual.}

\step
\recolor{sst.node["[0,7]"]}{state=current}
\narrate{Insert value 5 at position 2. Root $[0,7]$ partially covers -- recurse left.}

\step
\apply{sst}{add_node={id="[0,3]", parent="[0,7]"}}
\recolor{sst.node["[0,3]"]}{state=idle}
\narrate{Allocate left child $[0,3]$ on-demand. This node did not exist before.}

\step
\recolor{sst.node["[0,3]"]}{state=current}
\apply{sst}{add_node={id="[2,3]", parent="[0,3]"}}
\narrate{Recurse into $[0,3]$: position 2 is in the right half. Allocate $[2,3]$.}

\step
\recolor{sst.node["[2,3]"]}{state=current}
\apply{sst}{add_node={id="[2,2]", parent="[2,3]"}}
\narrate{Allocate leaf $[2,2]$ for position 2.}

\step
\apply{sst.node["[2,2]"]}{value=5}
\recolor{sst.node["[2,2]"]}{state=good}
\recolor{sst.node["[2,3]"]}{state=done}
\apply{sst.node["[2,3]"]}{value=5}
\recolor{sst.node["[0,3]"]}{state=done}
\apply{sst.node["[0,3]"]}{value=5}
\recolor{sst.node["[0,7]"]}{state=done}
\apply{sst.node["[0,7]"]}{value=5}
\narrate{Set leaf $[2,2] = 5$, bubble sums up: $[2,3]=5$, $[0,3]=5$, $[0,7]=5$.}

\step
\recolor{sst.node["[0,7]"]}{state=current}
\narrate{Insert value 3 at position 6. Root recurses right this time.}

\step
\apply{sst}{add_node={id="[4,7]", parent="[0,7]"}}
\recolor{sst.node["[4,7]"]}{state=idle}
\narrate{Allocate right child $[4,7]$.}

\step
\recolor{sst.node["[4,7]"]}{state=current}
\apply{sst}{add_node={id="[6,7]", parent="[4,7]"}}
\narrate{Position 6 is in $[6,7]$. Allocate $[6,7]$.}

\step
\recolor{sst.node["[6,7]"]}{state=current}
\apply{sst}{add_node={id="[6,6]", parent="[6,7]"}}
\narrate{Allocate leaf $[6,6]$.}

\step
\apply{sst.node["[6,6]"]}{value=3}
\recolor{sst.node["[6,6]"]}{state=good}
\recolor{sst.node["[6,7]"]}{state=done}
\apply{sst.node["[6,7]"]}{value=3}
\recolor{sst.node["[4,7]"]}{state=done}
\apply{sst.node["[4,7]"]}{value=3}
\recolor{sst.node["[0,7]"]}{state=done}
\apply{sst.node["[0,7]"]}{value=8}
\narrate{Leaf $[6,6]=3$. Bubble up: $[6,7]=3$, $[4,7]=3$, root $[0,7]=8$. Two paths allocated, rest of tree stays virtual.}

\step
\annotate{sst.node["[0,7]"]}{label="total = 8", color=good}
\narrate{Final state: only 8 nodes allocated out of a potential 15. Sparse trees use $O(q \log R)$ memory instead of $O(R)$.}
\end{animation}
